\chapter{Projet professionnel et perspectives de recherche}

\section{Projet de thèse : \enquote{Conception et Évaluation d'un Écosystème d'IA Co-adaptatif pour l'Inclusion Pédagogique Active}}

\subsection{Cadre et problématique}

Ce projet de thèse, envisagé au sein du laboratoire LITIS (UR 4108) sous la direction du Professeur Laurent Heutte et en collaboration avec Élisabeth SCHNEIDER (CIRNEF / ESO UMR 6590, Unicaen), vise à répondre à une problématique centrale de l'école inclusive : comment l'Intelligence Artificielle peut-elle devenir un partenaire co-adaptatif pour le triptyque élève-AESH-enseignant ?

L'objectif est de concevoir un écosystème socio-technique capable d'automatiser l'adaptation de contenus pédagogiques tout en maintenant une boucle de rétroaction humaine (Human-in-the-loop). Ce projet se situe à la confluence de trois domaines :
\begin{itemize}
    \item \textbf{Document Intelligence (IA)} : Pour l'analyse sémantique profonde des supports pédagogiques.
    \item \textbf{Interaction Homme-Machine (IHM)} : Pour garantir l'explicabilité et l'appropriation de l'outil par les enseignants (genèse instrumentale).
    \item \textbf{Apprentissage par Renforcement (RLHF)} : Pour aligner le système sur les objectifs pédagogiques réels via les retours des utilisateurs.
\end{itemize}

\subsection{Méthodologie et calendrier prévisionnel (2026-2029)}

La recherche se structurera en trois phases sur 3 ans :
\begin{enumerate}
    \item \textbf{Année 1 : Co-conception.} Immersion ethnographique et ateliers avec les utilisateurs (enseignants, AESH, élèves) pour modéliser les besoins et définir le cahier des charges.
    \item \textbf{Année 2 : Prototypage.} Développement d'un prototype open-source intégrant des modules d'adaptation multimodale (simplification de texte, description d'images, reformatage).
    \item \textbf{Année 3 : Expérimentation.} Déploiement en conditions réelles (IUT, lycée) et évaluation de l'impact sur l'autonomie des élèves et la charge de travail des accompagnants.
\end{enumerate}

Ce travail vise des contributions scientifiques (modélisation de l'IA co-adaptative), technologiques (outils open-source) et sociales (amélioration de l'inclusion).

\section{Trajectoire professionnelle et cohérence du projet VAE}

\subsection{Plan de carrière}

L'obtention du Master MEEF par VAE est la clé de voûte de ma transition professionnelle :
\begin{itemize}
    \item \textbf{Court terme (2025-2026)} : Consolidation de mon poste de PRCE à l'IUT du Havre et préparation de l'entrée en thèse.
    \item \textbf{Moyen terme (2026-2029)} : Réalisation du doctorat tout en poursuivant l'innovation pédagogique (Forge des Communs).
    \item \textbf{Long terme} : Qualification aux fonctions de Maître de Conférences pour diriger des recherches sur l'IA éducative.
\end{itemize}

\subsection{Le Master MEEF : un levier stratégique}

Ce diplôme ne constitue pas seulement une validation des acquis, mais un véritable tremplin. Il légitime 25 ans d'expérience technique et pédagogique, transformant un parcours atypique en une expertise académique reconnue. Il assure la continuité entre mon passé industriel, mon présent d'enseignant et mon futur de chercheur, avec pour fil rouge constant l'innovation au service de l'inclusion.
